
\SGARefTarget{sga3:xxiv:1.9}
\paragraph{\SGARefLink{sga3:xxiv:1.9}{Corollary 1.9}.}
\emph{Let \(S\) be a scheme, and let \(G\) and \(G'\) be two
reductive \(S\)-groups. Then}
\[
  F=\underline{\operatorname{Isom}}_{S\text{-}\mathrm{gr.}}(G,G')
\]
\emph{is representable by a smooth separated \(S\)-scheme, affine
if \(G\) or \(G'\) is semisimple. Moreover, \(S\) is the disjoint
union of two open subschemes \(S_1\) and \(S_2\) such that
\(F_{S_1}=\varnothing\), while \(F_{S_2}\) is a principal homogeneous
bundle on the left \textup{(respectively, on the right)} under
\(\underline{\Aut}_{S_2\text{-}\mathrm{gr.}}(G'_{S_2})\)
\textup{(respectively,
\(\underline{\Aut}_{S_2\text{-}\mathrm{gr.}}(G_{S_2})\)).}}

Indeed, let \(S_2\) be the set of points of \(S\) at which \(G\) and
\(G'\) have the same type, and let \(S_1\) be its complement. Since
the type of a reductive group is a locally constant function, \(S_1\)
and \(S_2\) are open. Clearly \(F_{S_1}=\varnothing\), and we may
assume \(S=S_2\). By \SGARefLink{sga3:XXIII:5.6}{Exp.~XXIII, 5.6}, \(F\) is a principal
homogeneous sheaf under
\(\underline{\Aut}_{S\text{-}\mathrm{gr.}}(G)\), locally trivial for
the étale topology. It follows that
\[
  F_0=F/\operatorname{ad}(G)
\]
is a principal homogeneous sheaf under \(\underline{\Autext}(G)\),
locally trivial for the étale topology, and hence representable
\textup{(\SGARefLink{sga3:x:corollary:5-5}{Exp.~X, 5.5})}. Since \(\operatorname{ad}(G)\) is affine,
\(F\) is therefore representable as well.

We note that the proof also gives the following result.

\SGARefTarget{sga3:xxiv:1.10}
\paragraph{\SGARefLink{sga3:xxiv:1.10}{Corollary 1.10}.}
\emph{Let \(S\) be a scheme, and let \(G\) and \(G'\) be two
reductive \(S\)-groups having the same type at every \(s\in S\).
Then \(\operatorname{ad}(G)\) acts freely on the right on
\(\underline{\operatorname{Isom}}_{S\text{-}\mathrm{gr.}}(G,G')\).
The quotient sheaf}
\[
  \underline{\operatorname{Isomext}}(G,G')
  =
  \underline{\operatorname{Isom}}_{S\text{-}\mathrm{gr.}}(G,G')
  /\operatorname{ad}(G)
\]
\emph{is representable by a twisted constant \(S\)-scheme; it is a
principal homogeneous bundle under \(\underline{\Autext}(G)\), and
is therefore finite over \(S\) if \(G\) is semisimple. Moreover, the
isomorphism}
\[
  \underline{\operatorname{Isom}}_{S\text{-}\mathrm{gr.}}(G,G')
  \simeq
  \underline{\operatorname{Isom}}_{S\text{-}\mathrm{gr.}}(G',G),
  \qquad u\longmapsto u^{-1},
\]
\emph{induces an isomorphism}
\[
  \underline{\operatorname{Isomext}}(G,G')
  \simeq
  \underline{\operatorname{Isomext}}(G',G).
\]

\SGARefTarget{sga3:xxiv:1.11}
\paragraph{\SGARefLink{sga3:xxiv:1.11}{Remark 1.11}.}
If
\(\underline{\operatorname{Isomext}}(G,G')(S)\neq\varnothing\), one
says that \(G\) is an \emph{inner twisted form} of \(G'\). Then \(G'\)
is an inner twisted form of \(G\), and the structure group of
\(\underline{\operatorname{Isom}}_{S\text{-}\mathrm{gr.}}(G,G')\)
can be reduced to \(\operatorname{ad}(G)\). More precisely, let
\[
  u\in\underline{\operatorname{Isomext}}(G,G')(S)
\]
and regard \(u\) as a section
\(u:S\to\underline{\operatorname{Isomext}}(G,G')\). Denote by
\[
  \underline{\operatorname{Isomint}}_u(G,G')
\]
the inverse image, under the canonical morphism
\[
  \underline{\operatorname{Isom}}_{S\text{-}\mathrm{gr.}}(G,G')
  \longrightarrow
  \underline{\operatorname{Isomext}}(G,G'),
\]
of the closed subscheme of
\(\underline{\operatorname{Isomext}}(G,G')\) that is the image of
\(u\). The natural action of \(\operatorname{ad}(G)\) on
\(\underline{\operatorname{Isomint}}_u(G,G')\) makes this scheme a
principal homogeneous bundle. By extension of the structure
group
\[
  \operatorname{ad}(G)\longrightarrow\underline{\Aut}(G),
\]
one recovers
\(\underline{\operatorname{Isom}}_{S\text{-}\mathrm{gr.}}(G,G')\)
from \(\underline{\operatorname{Isomint}}_u(G,G')\).

By Hensel's lemma \textup{(\SGARefLink{sga3:xi:corollary:1:11}{Exp.~XI, 1.11})}, \SGARefLink{sga3:xxiv:1.8}{Corollary~1.8} immediately
gives:
